\documentclass[12pt,a4paper]{article}
\usepackage[utf8]{inputenc}
\usepackage[T1]{fontenc}
\usepackage{amsmath,amssymb,amsthm}
\usepackage{listings}
\usepackage{xcolor}
\usepackage{graphicx}
\usepackage{hyperref}
\usepackage{geometry}
\usepackage{fancyhdr}
\usepackage{tcolorbox}
\usepackage{braket}

\geometry{margin=1in}

% Define colors for syntax highlighting
\definecolor{quantumblue}{RGB}{0,102,204}
\definecolor{classicalgreen}{RGB}{0,153,0}
\definecolor{governancered}{RGB}{204,0,0}
\definecolor{memoryorange}{RGB}{255,140,0}
\definecolor{commentgray}{RGB}{128,128,128}
\definecolor{backgroundgray}{RGB}{245,245,245}

% Configure listings for RiftLang pseudocode
\lstdefinestyle{riftlang}{
    backgroundcolor=\color{backgroundgray},
    basicstyle=\ttfamily\small,
    breaklines=true,
    captionpos=b,
    commentstyle=\color{commentgray}\itshape,
    keywordstyle=\color{quantumblue}\bfseries,
    keywordstyle=[2]\color{classicalgreen}\bfseries,
    keywordstyle=[3]\color{governancered}\bfseries,
    keywordstyle=[4]\color{memoryorange}\bfseries,
    morekeywords={@token, @operator, @memory_state, @gov_rule, @token_extension,
                  quantum, classical, QINT, INT, FLOAT, BRA, KET},
    morekeywords=[2]{collapse, superposition, entangle, measure},
    morekeywords=[3]{requires, prohibits, audit, enforce},
    morekeywords=[4]{lifetime, scope, alignment, persistence},
    numbers=left,
    numberstyle=\tiny\color{gray},
    frame=single,
    rulecolor=\color{black!30},
    showstringspaces=false,
    tabsize=2,
    xleftmargin=2em,
    framexleftmargin=1.5em,
    literate={|0>}{{$\ket{0}$}}3 {|1>}{{$\ket{1}$}}3 {|psi>}{{$\ket{\psi}$}}5
             {<psi|}{{$\bra{\psi}$}}5 {|q0>}{{$\ket{q_0}$}}4 {|q1>}{{$\ket{q_1}$}}4
             {|q2>}{{$\ket{q_2}$}}4 {|q3>}{{$\ket{q_3}$}}4 {|q4>}{{$\ket{q_4}$}}4
             {|q5>}{{$\ket{q_5}$}}4 {|q6>}{{$\ket{q_6}$}}4 {|q7>}{{$\ket{q_7}$}}4
             {hbar}{{$\hbar$}}1 {omega}{{$\omega$}}1 {pi}{{$\pi$}}1
             {<=}{{$\leq$}}1 {>=}{{$\geq$}}1 {->}{{$\rightarrow$}}1
             {alpha}{{$\alpha$}}1 {beta}{{$\beta$}}1 {epsilon}{{$\epsilon$}}1
             {tau}{{$\tau$}}1 {phi}{{$\phi$}}1 {psi}{{$\psi$}}1
             {lambda}{{$\lambda$}}1 {Gamma}{{$\Gamma$}}1 {Delta}{{$\Delta$}}1
             {^2}{{$^2$}}1 {^-44}{{$^{-44}$}}3 {^-15}{{$^{-15}$}}3 {^-9}{{$^{-9}$}}2
             {_2}{{$_2$}}1 {10^6}{{$10^6$}}3 {10^-44}{{$10^{-44}$}}5
}

% Header configuration
\pagestyle{fancy}
\fancyhf{}
\lhead{RiftLang Protocol Stack - Stage-N Specification}
\rhead{AEGIS Compliance Framework}
\cfoot{\thepage}

\title{
    \vspace{-2cm}
    \Large \textbf{AEGIS Project} \\
    \vspace{0.5cm}
    \Huge \textbf{Stage-N: Hybrid Quantum-Classical Token Execution Layer} \\
    \vspace{0.5cm}
    \Large \textbf{Formal Specification for RiftLang Protocol Stack} \\
    \vspace{0.5cm}
    \normalsize Version 1.0 - Technical Specification Document
}

\author{
    OBINexus Computing Division \\
    \texttt{Toolchain: riftlang.exe $\rightarrow$ .so.a $\rightarrow$ rift.exe $\rightarrow$ gosilang}
}

\date{\today}

\begin{document}

\maketitle
\thispagestyle{empty}

\newpage
\tableofcontents
\newpage

\section{Executive Summary}

This document formalizes Stage-N of the RiftLang Protocol Stack, establishing the critical interface between quantum probabilistic computation and classical deterministic execution. Stage-N enables seamless transitions between quantum superposition states and classical computational models while maintaining strict AEGIS governance compliance throughout the quantum-classical boundary.

The specification defines standardized patterns for stages 0 through N+1 (currently implemented through Stage-7), providing a unified framework for quantum token management, collapse operations, and memory-governed parsing within the RIFT domain-specific language ecosystem.

\section{Stage Evolution Framework}

\begin{tcolorbox}[colback=blue!5!white,colframe=quantumblue,title=Stage Progression Model]
\begin{itemize}
    \item \textbf{Stage-0}: Token initialization and classical baseline
    \item \textbf{Stage-1}: Quantum extension introduction
    \item \textbf{Stage-2}: Entanglement protocol establishment
    \item \textbf{Stage-3}: Collapse operator implementation
    \item \textbf{Stage-4}: Memory governance integration
    \item \textbf{Stage-5}: Parser unification
    \item \textbf{Stage-6}: AEGIS phase alignment
    \item \textbf{Stage-7}: Full quantum-classical bridge deployment
    \item \textbf{Stage-N}: Dynamic stage instantiation
    \item \textbf{Stage-N+1}: Future extensibility framework
\end{itemize}
\end{tcolorbox}

\section{Core Architecture and Purpose}

\subsection{Fundamental Design Principles}

Stage-N serves as the quantum-classical computation bridge within the RIFT DSL execution pipeline. The architecture enables:

\begin{enumerate}
    \item \textcolor{quantumblue}{Quantum state preservation} during computation
    \item \textcolor{classicalgreen}{Deterministic resolution} when measurement occurs
    \item \textcolor{governancered}{Governance enforcement} at state transitions
    \item \textcolor{memoryorange}{Memory-bounded execution} with Planck-scale constraints
\end{enumerate}

\subsection{Integration Points}

\begin{lstlisting}[style=riftlang,caption={Stage-N Integration Architecture}]
@stage_interface[N] {
    input_stages: [N-1, quantum_init]
    output_stages: [N+1, classical_exec]
    governance_hooks: AEGIS_Phase_I_III
    memory_model: quantum_foam_temporal
    compliance_level: STRICT
}
\end{lstlisting}

\section{Quantum Token Specification}

\subsection{Core Quantum Token Definition}

The fundamental quantum token represents a qubit in superposition state with complex amplitude coefficients.

\begin{lstlisting}[style=riftlang,caption={Quantum Token Base Definition}]
@token[quantum] qbit superposition(alpha, beta) -> QINT
Where:
    alpha in C : Complex amplitude for |0> state
    beta  in C : Complex amplitude for |1> state
    Constraint: |alpha|^2 + |beta|^2 = 1 (normalization)
\end{lstlisting}

\subsection{Extended Quantum Token Attributes}

\begin{lstlisting}[style=riftlang,caption={Quantum Token Extensions}]
@token_extension[quantum] {
    # Dirac notation representation
    bra_ket_notation: |psi> = alpha|0> + beta|1>
    
    # Normalization enforcement with tolerance
    amplitude_norm: enforce(|alpha|^2 + |beta|^2 = 1.0 +/- epsilon)
    Where: epsilon = 10^-15 (machine epsilon)
    
    # Decoherence threshold at Planck scale
    decohere_threshold: tau_planck = 5.39 x 10^-44 seconds
    
    # Entanglement tracking
    entanglement_flag: bool
    entanglement_partners: QINT[] (max_size = 6)
    
    # Phase coherence bounds
    phase_coherence: phi in [0, 2*pi]
    phase_drift_rate: d_phi/dt <= pi/tau_coherence
}
\end{lstlisting}

\subsection{Quantum Token Memory Alignment}

\begin{lstlisting}[style=riftlang,caption={Memory Alignment Specification}]
@memory_align[quantum] {
    alignment: 8-qubit boundary
    span_type: distributed_quantum_foam
    coherence_window: planck_time
    isolation_level: phase_locked
    
    layout: {
        |q0>|q1>|q2>|q3>|q4>|q5>|q6>|q7>
        [amplitude_real][amplitude_imag]
        [phase][entangle_mask][coherence]
    }
}
\end{lstlisting}

\section{Classical Resolution Operator}

\subsection{Collapse Operator Definition}

The collapse operator manages the quantum-to-classical transition under governance constraints.

\begin{lstlisting}[style=riftlang,caption={Collapse Operator Specification}]
@operator collapse {
    input: QINT                    # Quantum integer in superposition
    condition: coherence >= PLANCK_THRESHOLD
    output: INT                    # Classical deterministic integer
    audit: quantum_event_log       # Governance audit trail
    
    properties: {
        irreversible: true
        measurement_basis: computational
        entropy_increase: Delta_S > 0
    }
}
\end{lstlisting}

\subsection{Piecewise Collapse Logic}

\begin{lstlisting}[style=riftlang,caption={Collapse Decision Tree}]
PROCEDURE quantum_classical_collapse(q: QINT, E: energy, m: mass):
    LET c^2 = speed_of_light^2
    LET E_planck = hbar * omega
    
    IF E^2 <= m * c^2:
        # Low energy - classical behavior dominates
        state := CLASSICAL_DETERMINISTIC
        RETURN cast_to_int(q, method="measurement")
        
    ELIF E > E_critical AND q.coherence >= tau_planck:
        # High energy with coherence - forced collapse
        TRIGGER collapse_event {
            log: quantum_event_log
            timestamp: current_planck_time
            method: "forced_decoherence"
        }
        state := CLASSICAL_COLLAPSED
        RETURN probabilistic_cast(q)
        
    ELSE:
        # Maintain quantum superposition
        state := QUANTUM_SUPERPOSITION
        EVOLVE q WITH hamiltonian(H)
        RETURN q  # Preserve quantum state
        
    END IF
END PROCEDURE
\end{lstlisting}

\subsection{State Transition Matrix}

\begin{lstlisting}[style=riftlang,caption={Quantum-Classical State Transitions}]
@state_transition_matrix {
    QUANTUM -> CLASSICAL: {
        trigger: measurement OR decoherence
        probability: |<psi|phi>|^2
        governance: collapse_contract
        audit_level: MANDATORY
    }
    
    CLASSICAL -> QUANTUM: {
        trigger: superposition_gate
        condition: coherence_budget > threshold
        governance: quantum_init_contract
        audit_level: STRICT
    }
    
    QUANTUM -> QUANTUM: {
        trigger: unitary_evolution
        operator: U = exp(-iHt/hbar)
        governance: evolution_contract
        audit_level: PERIODIC
    }
}
\end{lstlisting}

\section{Memory-Governed Quantum Parser}

\subsection{Hybrid Token Grammar}

The parser must handle both quantum and classical tokens with appropriate type safety.

\begin{lstlisting}[style=riftlang,caption={Token Type Definitions}]
@parser[hybrid_quantum_classical] {
    token_types: {
        QINT    : quantum_integer[superposition]
        INT     : classical_integer[deterministic]
        FLOAT   : classical_float[ieee754]
        BRA     : <psi| quantum_state
        KET     : |psi> quantum_state
        QFLOAT  : quantum_float[superposition]
    }
    
    # Regex pattern with quantum extensions
    parse_rules: R"/([QC])(INT|FLOAT|STATE)/gmi[tb]"
    Where:
        g: global matching
        m: multiline quantum states
        i: case-insensitive operators
        t: top-down classical resolution
        b: bottom-up quantum composition
}
\end{lstlisting}

\subsection{Temporal Memory State Management}

\begin{lstlisting}[style=riftlang,caption={Quantum Memory Management}]
@memory_state::quantum_foam {
    lifetime: planck_time = 5.39 x 10^-44 seconds
    scope: local_superposition
    
    allocation: {
        classical_mode: align(4096_bits)
        quantum_mode: align(8_qubits)
        hybrid_mode: interleaved_coherent
    }
    
    persistence: {
        coherent_duration: tau_coherence
        decoherence_rate: Gamma = 1/T_2
        error_threshold: 10^-9
    }
    
    governance: {
        max_entanglement_depth: 6
        bell_state_limit: 4_pairs
        gc_policy: phase_aware_collection
    }
}
\end{lstlisting}

\subsection{Parser State Machine}

\begin{lstlisting}[style=riftlang,caption={Quantum Parser Automaton}]
AUTOMATON quantum_parser {
    states: {S_INIT, S_QUANTUM, S_CLASSICAL, S_COLLAPSE, S_MEASURE}
    
    transitions: {
        S_INIT -> S_QUANTUM: 
            condition: detect_superposition_token()
            action: init_quantum_context()
            
        S_QUANTUM -> S_COLLAPSE:
            condition: coherence < PLANCK_THRESHOLD
            action: prepare_collapse()
            
        S_COLLAPSE -> S_CLASSICAL:
            condition: collapse_complete()
            action: emit_classical_token()
            
        S_QUANTUM -> S_MEASURE:
            condition: measurement_operator()
            action: von_neumann_projection()
            
        S_MEASURE -> S_CLASSICAL:
            condition: measurement_complete()
            action: emit_measured_value()
    }
    
    error_states: {
        E_COHERENCE_LOST: recovery = forced_collapse
        E_ENTANGLE_VIOLATION: recovery = isolate_subsystem
        E_MEMORY_OVERFLOW: recovery = quantum_gc
    }
}
\end{lstlisting}

\section{Governance Constraint Declarations}

\subsection{Core Governance Rules}

\begin{lstlisting}[style=riftlang,caption={Collapse Contract Governance}]
@gov_rule::collapse_contract {
    requires: {
        coherence >= planck_threshold
        entanglement_depth <= max_allowed
        quantum_budget > operation_cost
        audit_trail.enabled = true
    }
    
    prohibits: {
        superposition_state > max_density_matrix_size
        concurrent_measurements > 1
        phase_drift > pi/4
        untracked_entanglement = true
    }
    
    audit: {
        log_destination: quantum_event_log
        retention_period: 7_stages
        cryptographic_seal: SHA3-256
        immutability: blockchain_anchored
    }
}
\end{lstlisting}

\subsection{Resource Management Constraints}

\begin{lstlisting}[style=riftlang,caption={Quantum Resource Governance}]
@gov_rule::quantum_resource_management {
    allocation_policy: {
        max_qubits_per_token: 16
        max_entangled_pairs: 8
        decoherence_budget: 1000_planck_times
        memory_quota: 1MB_quantum_foam
    }
    
    cleanup_policy: {
        auto_collapse_timeout: 100_planck_times
        garbage_collection: phase_aware
        memory_reclaim: immediate
        entanglement_pruning: depth_first
    }
    
    cost_model: {
        superposition_cost: 0.1_per_qubit_per_cycle
        entanglement_cost: 0.3_per_pair
        measurement_cost: 0.2_per_operation
        coherence_maintenance: 0.05_per_planck_time
    }
}
\end{lstlisting}

\subsection{Security and Validation Rules}

\begin{lstlisting}[style=riftlang,caption={Quantum Security Governance}]
@gov_rule::quantum_security {
    validation: {
        state_vector_normalization: continuous
        no_cloning_enforcement: strict
        basis_state_verification: periodic(10_cycles)
        bell_inequality_check: on_entanglement
    }
    
    access_control: {
        quantum_state_read: privileged_only
        collapse_trigger: authorized_operators
        entanglement_create: rate_limited(10/sec)
        phase_manipulation: governance_approved
    }
    
    integrity: {
        checksum_algorithm: quantum_hash_SHA3Q
        tamper_detection: bell_inequality_test
        audit_trail: immutable_quantum_ledger
        replay_protection: nonce_per_operation
    }
}
\end{lstlisting}

\section{Integration with AEGIS Phase Architecture}

\subsection{Phase I - Matrix Parity Integration}

\begin{lstlisting}[style=riftlang,caption={Matrix Parity Bridge}]
INTEGRATION matrix_parity_bridge {
    quantum_to_fft: {
        INPUT: QINT[superposition]
        PROCESS: 
            1. Extract amplitude vectors (alpha, beta)
            2. Map to FFT basis: F(|psi>) = Sum(alpha_k*e^(2*pi*ijk/N))
            3. Apply parity constraints from Phase I
            4. Verify matrix eigenvalue stability
        OUTPUT: FFT_MATRIX[classical]
    }
    
    governance: {
        parity_check: R"/[01]{8}/g" 
        matrix_alignment: 8x8_quantum_block
        eigenvalue_threshold: |lambda| < 1.0
    }
}
\end{lstlisting}

\subsection{Phase II - Token Stream Management}

\begin{lstlisting}[style=riftlang,caption={Token Stream Integration}]
INTEGRATION token_stream_quantum {
    stream_mode: {
        classical: sequential_ordered
        quantum: parallel_superposed
        hybrid: context_switched
    }
    
    synchronization: {
        barrier: quantum_measurement_point
        ordering: causal_cone_preservation
        latency: <= coherence_window
    }
    
    buffering: {
        quantum_buffer: circular_phase_locked
        classical_buffer: FIFO_deterministic
        transition_buffer: copy_on_collapse
    }
}
\end{lstlisting}

\subsection{Phase III - Planck Verification}

\begin{lstlisting}[style=riftlang,caption={Planck Scale Verification}]
INTEGRATION planck_verification {
    collapse_window: {
        detection: coherence < PLANCK_THRESHOLD
        action: enforce(collapse_contract)
        verification: cryptographic_proof
        timing: exact_planck_time
    }
    
    quantum_classical_boundary: {
        transition_log: {
            timestamp: planck_time_resolution
            state_before: |psi>
            state_after: classical_value
            entropy_change: Delta_S
            information_preserved: I = -Sum(p log p)
        }
    }
    
    entanglement_boundary: {
        max_distance: 6_hops
        isolation_enforcement: bell_state_collapse
        audit_requirement: full_trace
        correlation_preservation: EPR_compliant
    }
}
\end{lstlisting}

\section{Runtime Execution Model}

\subsection{Dual-Mode Execution Pipeline}

\begin{lstlisting}[style=riftlang,caption={Stage-N Execution Pipeline}]
PIPELINE rift_stage_n_execution {
    MODE classical {
        stages: tokenize -> parse -> validate -> execute
        memory: sequential_4096_aligned
        concurrency: mutex_protected
        error_handling: exception_based
    }
    
    MODE quantum {
        stages: superpose -> entangle -> evolve -> measure
        memory: distributed_8qubit_foam
        concurrency: phase_locked_parallel
        error_handling: decoherence_recovery
    }
    
    MODE hybrid {
        stages: detect_context -> switch_mode -> process -> reconcile
        memory: adaptive_alignment
        concurrency: quantum_classical_barrier
        error_handling: graceful_degradation
    }
    
    performance: {
        classical_throughput: 10^6 ops/sec
        quantum_coherence: 1000 x planck_time
        transition_overhead: < 1us
    }
}
\end{lstlisting}

\subsection{Context Switching Protocol}

\begin{lstlisting}[style=riftlang,caption={Quantum-Classical Context Switch}]
PROTOCOL context_switch {
    classical_to_quantum: {
        save_classical_state()
        init_quantum_registers()
        prepare_superposition()
        verify_coherence()
        enable_quantum_operations()
    }
    
    quantum_to_classical: {
        prepare_measurement()
        collapse_superposition()
        extract_classical_value()
        cleanup_quantum_resources()
        restore_classical_context()
    }
    
    switch_overhead: {
        time_cost: O(n_qubits)
        memory_cost: O(2^n_qubits)
        governance_cost: O(audit_depth)
    }
}
\end{lstlisting}

\section{Example Usage Patterns}

\subsection{Basic Quantum Token Operations}

\begin{lstlisting}[style=riftlang,caption={Quantum Token Usage Examples}]
# Initialize quantum token in superposition
@quantum
let q_value: QINT = superposition(0.707, 0.707)  # |+> state

# Entangle two quantum tokens
@quantum
let q_pair: (QINT, QINT) = entangle(q_value, q_other)

# Conditional collapse based on coherence
@hybrid
if coherence(q_value) < PLANCK_THRESHOLD {
    let classical_result: INT = collapse(q_value)
    process_classical(classical_result)
} else {
    evolve_quantum(q_value, hamiltonian)
}

# Measurement with basis selection
@quantum
let measured: INT = measure(q_value, basis="X")
\end{lstlisting}

\subsection{Hybrid Computation Pattern}

\begin{lstlisting}[style=riftlang,caption={Hybrid Quantum-Classical Algorithm}]
@hybrid_algorithm grover_search {
    # Classical preprocessing
    let dataset: INT[] = load_classical_data()
    let target: INT = define_search_target()
    
    # Quantum acceleration phase
    @quantum {
        let q_register: QINT[] = init_superposition(size=log2(dataset.length))
        
        repeat sqrt(dataset.length) times {
            apply_oracle(q_register, target)
            apply_diffusion(q_register)
            
            if coherence_degraded(q_register) {
                refresh_quantum_state(q_register)
            }
        }
        
        let result_index: INT = measure_all(q_register)
    }
    
    # Classical verification
    @classical {
        verify_result(dataset[result_index], target)
        return dataset[result_index]
    }
}
\end{lstlisting}

\section{Formal Verification Properties}

\subsection{Safety Properties}

\begin{lstlisting}[style=riftlang,caption={Quantum Safety Invariants}]
PROPERTY quantum_safety {
    # G1: Normalization is always maintained
    [] (forall q: QINT . |amplitude(q)|^2 = 1.0 +/- epsilon)
    
    # G2: No-cloning theorem is preserved
    [] (not exists operation . clone(|psi>) -> |psi>|psi>)
    
    # G3: Causality is respected
    [] (forall measurement . timestamp(cause) < timestamp(effect))
    
    # G4: Measurement irreversibility
    [] (collapsed(q) -> not quantum(q))
}
\end{lstlisting}

\subsection{Liveness Properties}

\begin{lstlisting}[style=riftlang,caption={Quantum Liveness Guarantees}]
PROPERTY quantum_liveness {
    # L1: Every quantum state eventually decoheres
    <> (forall q: QINT . coherence(q) < PLANCK_THRESHOLD)
    
    # L2: Measurements eventually complete
    (measure_initiated(q) -> <> measure_complete(q))
    
    # L3: Resources are eventually reclaimed
    []<> (allocated_qubits = deallocated_qubits)
    
    # L4: Entanglements eventually resolve
    <> (forall entangled_pair . separated OR measured)
}
\end{lstlisting}

\section{Performance Specifications}

\begin{tcolorbox}[colback=green!5!white,colframe=classicalgreen,title=Performance Requirements]
\begin{itemize}
    \item \textbf{Quantum Coherence Time}: $\geq 1000 \tau_{planck}$
    \item \textbf{State Preparation}: $< 10$ ns per qubit
    \item \textbf{Measurement Latency}: $< 100$ ns
    \item \textbf{Context Switch Overhead}: $< 1 \mu$s
    \item \textbf{Memory Efficiency}: $O(n)$ for $n$ qubits
    \item \textbf{Entanglement Depth}: Maximum 6 levels
    \item \textbf{Error Rate}: $< 10^{-9}$ per operation
\end{itemize}
\end{tcolorbox}

\section{Firmware Integration Guidelines}

\subsection{Git-RAF Integration Points}

\begin{lstlisting}[style=riftlang,caption={Git-RAF Firmware Hooks}]
@firmware_integration {
    git_raf_hooks: {
        pre_commit: validate_quantum_governance()
        post_merge: verify_collapse_consistency()
        pre_push: audit_quantum_operations()
    }
    
    versioning: {
        quantum_state_snapshot: on_commit
        collapse_history: immutable_log
        entanglement_graph: version_tracked
    }
    
    deployment: {
        stage: [0..N+1]
        firmware_target: quantum_coprocessor
        validation_level: AEGIS_COMPLIANT
    }
}
\end{lstlisting}

\section{Conclusion and Future Extensions}

Stage-N of the RiftLang Protocol Stack provides a robust foundation for hybrid quantum-classical computation within the AEGIS governance framework. The specification enables:

\begin{itemize}
    \item Seamless quantum-classical transitions
    \item Governance-compliant state management
    \item Planck-scale temporal constraints
    \item Integration with existing AEGIS phases
    \item Extensibility for future quantum algorithms
\end{itemize}

Future stages (N+1 and beyond) will extend this framework to support:
\begin{itemize}
    \item Distributed quantum computation
    \item Fault-tolerant quantum error correction
    \item Advanced entanglement protocols
    \item Quantum machine learning integration
\end{itemize}

\appendix

\section{Glossary of Terms}

\begin{description}
    \item[QINT] Quantum Integer - A quantum register holding superposition states
    \item[Planck Time] $\tau_{planck} = 5.39 \times 10^{-44}$ seconds
    \item[Coherence] Quantum state phase relationship preservation
    \item[Collapse] Quantum to classical state transition
    \item[AEGIS] Automated Enterprise Governance Intelligence System
\end{description}

\section{References}

\begin{enumerate}
    \item AEGIS Project Technical Specification v2.0
    \item RiftLang Compiler Documentation
    \item Quantum Computing Governance Framework
    \item Git-RAF Firmware Integration Manual
    \item OBINexus Toolchain Architecture Guide
\end{enumerate}

\end{document}